\documentclass[10pt,a4paper]{article}
\usepackage[landscape,margin=14mm,top=22mm,bottom=18mm,headheight=25pt]{geometry}
\usepackage[utf8]{inputenc}
\usepackage[T1]{fontenc}
\usepackage[brazil]{babel}
\usepackage{helvet,amsmath,amssymb,booktabs,array,tabularx,enumitem,graphicx,xcolor,tcolorbox,fancyhdr,hyperref,tikz}
\usetikzlibrary{arrows.meta,positioning}
\renewcommand{\familydefault}{\sfdefault}
\definecolor{azul}{RGB}{0,83,134}
\definecolor{ouro}{RGB}{206,177,68}
\hypersetup{pdflang=pt-BR,colorlinks=true,urlcolor=azul,linkcolor=azul,pdfauthor={Luan F. Nogueira e Luiz F. de Melo Soares}}
\setlength{\parindent}{0pt}
\setlength{\parskip}{5pt}
\setlist{nosep,leftmargin=*}
\pagestyle{fancy}\fancyhf{}
\fancyhead[L]{\small\color{azul}DO TRANSISTOR AO CHIP}
\fancyhead[R]{\includegraphics[height=7mm]{../slides/libs/emblemufc.pdf}}
\fancyfoot[L]{\scriptsize Luan F. Nogueira • Luiz F. de Melo Soares}
\fancyfoot[R]{\scriptsize\thepage}
\renewcommand{\headrulewidth}{.4pt}
\setlength{\headheight}{25pt}
\newcommand{\titulo}[2]{{\LARGE\color{azul}\bfseries #1}\par{\large #2}\par\vspace{2mm}}
\newcommand{\caixa}[2]{\begin{tcolorbox}[colback=azul!4,colframe=azul,title={#1},fonttitle=\bfseries,boxrule=.5pt,arc=1mm]#2\end{tcolorbox}}
\newcommand{\VDD}{V_{\mathrm{DD}}}

\usepackage{pgfplots}\pgfplotsset{compat=1.18}
\definecolor{blue_theme}{RGB}{0,83,134}\definecolor{ativo}{RGB}{0,132,75}\definecolor{lacuna}{RGB}{199,123,0}
\makeatletter\def\input@path{{../slides/}}\makeatother
\input{integracao_visual.tex}
\begin{document}
\titulo{Histórico de integração em processadores}{Luan F. Nogueira e Luiz F. de Melo Soares · contagens por die · seleção 1971--2024}
\begin{minipage}[t]{.48\linewidth}\vspace{0pt}\centering\resizebox{\linewidth}{!}{\graficoIntegracao{2}}\par\smallskip\footnotesize O gráfico amplia o recorte 2020--2024; a tabela reúne a série completa. CPU, SoC e GPU têm funções e áreas diferentes. Contagem não equivale a desempenho ou densidade. Nós de processo não medem uniformemente o comprimento do Gate.\end{minipage}\hfill
\begin{minipage}[t]{.49\linewidth}\vspace{0pt}\scriptsize\begin{tabular}{rllrr}\toprule Ano&Chip&Classe&Transistores&Processo\\\midrule
1971&4004&CPU&2,3 mil&10 $\mu$m\\
1978&8086&CPU&29 mil&3 $\mu$m\\
1985&386 DX&CPU&275 mil&1,5 $\mu$m\\
1993&Pentium&CPU&3,1 mi&800 nm\\
2000&Pentium 4&CPU&42 mi&180 nm\\
2006&Core 2 Duo&CPU&291 mi&65 nm\\
2008&Core i7&CPU&731 mi&45 nm\\
2012&Ivy Bridge&CPU + GPU&1,4 bi&22 nm\\
2020&M1&SoC Arm&16 bi&5 nm\\
2020&A100&GPU&54,2 bi&7 nm N7\\
2022&M2&SoC Arm&20 bi&5 nm (2ª ger.)\\
2023&M3&SoC Arm&25 bi&3 nm\\
2024&M4&SoC Arm&28 bi&3 nm (2ª ger.)\\
\bottomrule\end{tabular}\par\medskip\footnotesize Fontes: Intel, Apple, NVIDIA e artigo de Ivy Bridge no ISSCC 2012. Links individuais, revisões selecionadas e observações no CSV e no documento FONTES\_INTEGRACAO.md.\end{minipage}
\end{document}
